% Paper outline for review, per Allison's format:
%   - title
%   - authors
%   - complete abstract (placeholders OK for unknown numbers)
%   - outline down to subsection level
%   - all expected figures/tables and captions (placeholder/sketch OK)
% This is a separate deliverable from the thesis (main.tex / sections/).
% It mirrors the thesis structure but is condensed for review, and
% marks open items as [TODO] / PLACEHOLDER rather than writing full
% prose, the point is to align on scope and structure before the
% thesis content is finalized.
\documentclass[11pt]{article}
\usepackage[utf8]{inputenc}
\usepackage[margin=1in]{geometry}
\usepackage{graphicx}
\usepackage{amsmath}
\usepackage{cite}
\usepackage{hyperref}

\title{Real-Time Multimodal Performance Assessment for dVRK \\
\large Paper Outline: Draft for Review}

% TODO (with the group): confirm author order and affiliations.
% Placeholder order below: student first, mentors, co-supervisor,
% supervisor/corresponding last, standard convention, not yet
% confirmed with anyone listed.
\author{
  Shujiro Shobayashi$^{1,2}$ \and
  Brian Vuong$^{1}$ \and % TODO: confirm affiliation
  Mary Kate Gale$^{1}$ \and % TODO: confirm affiliation
  Josephine Hughes$^2$ \and
  Allison M. Okamura$^{1}$ \\[2mm]
  \normalsize $^1$Stanford University \quad $^2$EPFL
}
\date{\today}

\begin{document}
\maketitle

\section*{Abstract}
% Complete structurally; numeric results are placeholders pending data.
Real-time performance feedback can accelerate skill acquisition during
robotic surgical training, but existing systems either compute
performance measures only after a task ends (post-hoc), or depend on a human
rater to trigger and score feedback during the trial. We present a
multimodal system that measures peg-transfer task performance live, on
a physical da Vinci Research Kit, with no human rater at any
point in the loop. The system fuses kinematic, visual, and force data
through an event-triggered pipeline: trial time and instrument path
length are tracked in real time, cylinder orientation is checked at
placement via HSV segmentation on a YOLOv8-detected
cylinder, and contact force is reported immediately after each trial.
In addition, with the depth camera, we collect a higher-accuracy,
SAM2-based 3D localization measure post-hoc, as extra
validation data for the
real-time measures. We demonstrate the system on a physical dVRK, validating the real-time measures against a higher-accuracy post-hoc reference, and show that automated, real-time, on-robot performance measurement is feasible for a structured surgical training task; laying the groundwork for a closed-loop, difficulty-adaptive training system.

\section*{Outline}

\subsection*{1. Introduction}
\begin{itemize}
  \item \textbf{1.1 Scope}: Measures performance on one well-defined
  task (peg transfer), not global surgical skill; task-level metrics
  are objective and real-time-computable, a global skill judgment is
  neither.
  \item \textbf{1.2 Goal}: Real-time dVRK peg-transfer monitoring,
  computed directly from sensor/vision data, no human rater. Five
  measures tracked: completion time, path length, rate of orientation
  change, angular displacement, object state. Long-term vision:
  closed-loop, difficulty-adaptive task.
  \item \textbf{1.3 Motivation}: (a) real-time feedback accelerates
  learning~\cite{laca2022}; (b) automated capture catches events a
  human observer misses; (c) multimodal data can surface metrics beyond
  completion time.
  \item \textbf{1.4 Contribution}: Closes the gap between offline
  accurate measurement and real-time-but-human-scored feedback: a
  multimodal, automated system, live, on the physical robot, for one
  task.
\end{itemize}

\subsection*{2. Related Work}
\begin{itemize}
  \item \textbf{2.1 Post-Hoc Analysis}: JIGSAWS-style~\cite{jigsaws}
  kinematics and video, both offline: skill proxies (path length,
  velocity, jerk) computed from recorded kinematics; object detection
  on recorded video for task-progress tracking, accurate but too heavy
  for real-time.
  \item \textbf{2.2 Real-Time Multimodal Skill Prediction}: Closest
  technical precedent (Hu et al.~\cite{hu2026}): kinematic + vision
  fusion, real-time-capable model, but evaluated offline on recorded
  trials and predicts a general OSATS score, not task-specific metrics.
  \item \textbf{2.3 Real-Time Feedback Attempts}: Closest in spirit
  (Laca et al.~\cite{laca2022}): real-time feedback RCT on physical
  hardware, but the performance measure behind the feedback is scored
  post-hoc by blinded human raters.
\end{itemize}


\begin{table}[h]
\centering
\begin{tabular}{|p{3.3cm}|p{5.1cm}|p{4.3cm}|}
\hline
\textbf{Approach} & \textbf{What it does} & \textbf{Limitation} \\
\hline
Post-hoc analysis (JIGSAWS-style: kinematics + video) &
Path length, velocity, jerk from kinematics; object detection from
video for task-progress &
Offline; not live; vision side also computationally heavy \\[4pt]
\hline
Real-time multimodal skill prediction (Hu et al.~\cite{hu2026}) &
Fuses kinematic + vision data from the da Vinci system to predict OSATS
skill score in real time &
Evaluated offline on recorded trials; predicts general skill, not
task-specific metrics \\[4pt]
\hline
Real-time feedback (Laca et al.~\cite{laca2022}) &
RCT: real-time feedback during a robotic dissection task, on physical
hardware &
Performance is still scored post-hoc, by blinded human video review \\[4pt]
\hline
\textbf{This work} &
Multimodal, automated, on the physical robot, live &
None (the gap this thesis fills) \\
\hline
\end{tabular}
\caption{Related work on surgical task assessment, compared against the
requirements this thesis targets, matching \texttt{tab:related\_work}
in \texttt{sections/background.tex}.}
\label{tab:outline_related_work}
\end{table}


\subsection*{3. System \& Task Setup}
\begin{itemize}
  \item \textbf{3.1 Hardware}: dVRK: da Vinci S console, da Vinci Si
  patient-side arms, PSM1/PSM2, with bipolar fenestrated forceps.
  \item \textbf{3.2 Formative Study}: Prior data collection (27
  novices, 9 experts, robotic surgeons with $>$250 hours of
  laparoscopic teleoperated experience) for Mary Kate Gale's study on
  fine wrist control as a marker of teleoperation expertise; the author
  assisted in data collection. Findings: participant-preferred camera
  views; frustration with unimanual control (PSM1 only, no clutch, no
  camera repositioning); with novice operators, non-smooth motion that
  damaged pegs on the board. High success rates plus peg breakage
  motivated a new, lower-peg-height board, built but not used in this
  thesis; instead it, and these findings, motivate the closed-loop,
  difficulty-adaptive task proposed as future work (Section 7).
  \item \textbf{3.3 Task}: Cylinder-on-peg transfer: 4 center pegs
  (rest position), 8 outer pegs with blue/white LEDs (target
  indication), conductive-paint board for ground-truth contact
  detection. \emph{[Figure 1: task board]}
  \item \textbf{3.4 Vision}: dVRK endoscopic stereo camera; Intel
  RealSense D405 depth camera, fixed external view; External camera for additional view, pedal-triggered.
  \item \textbf{3.5 Sensing}: PSM kinematics; Arduino-based
  task start/end trigger; ATI force sensor under the task pad.
\end{itemize}

\subsection*{4. Methods}
\begin{itemize}
  \item \textbf{4.1 Pipeline Overview}: One pipeline, built from
  modules the user can enable independently at session start (trial
  trigger, path length, orientation rate, force sensor, orientation
  check). Walkthrough with everything on: lift $\to$ Arduino signal
  $\to$ clock starts; place $\to$ Arduino signal $\to$ clock stops
  $\to$ HSV orientation check; end of trial $\to$ force popup (peak,
  threshold crossings, graph). \emph{[Figure 2: pipeline diagram;
  Table 2: data/timing/feedback summary]}
  \item \textbf{4.2 Real-Time GUI \& Trial Clock}: Trigger is a
  session choice: MONO pedal + teleop (clock/path length run only
  while pedal held) or Arduino cylinder lift/place (clock runs on its
  own); this thesis's results use Arduino. 200~Hz clock; two
  always-on-top windows; color-coded progress bar; live path length,
  both arms. \emph{[Figure 3: GUI screenshot]} [TODO: source code
  reference/appendix]
  \item \textbf{4.3 Kinematic Metrics}: Path length and rate of
  orientation change, both computed directly from PSM kinematics.
  Path length: cumulative Euclidean distance between tool-tip
  samples. Rate of orientation change: quaternion difference between
  orientation samples $\to$ rotation angle $\to$ divided by elapsed
  time $\to$ averaged over the trial, following Equation~8 of Sharon
  et al.~\cite{sharon2021}, who found this metric separates expert
  from novice performance alongside task time and path length.
  (Formulas in thesis body.)
  \item \textbf{4.4 Ground Contact Detection}: ATI sensor on a
  separate computer, streamed over UDP to the pipeline; read live
  (real-time capable), smoothed with a 5-sample moving average.
  Shown to the trainee as a post-trial popup instead, by design, so
  it doesn't compete for attention with the clock/path length during
  the trial: peak (smoothed) force, count of crossings over a 5~N
  threshold, and a force/time graph, all from the same smoothed
  signal. \emph{[Figure 4: force graph]}
  \item \textbf{4.5 Object Detection}: YOLOv8s~\cite{yolov8}, 4 classes
  (cylinder, inactive/blue/white peg). Trained on frames from the
  task's own recordings, labeled on Roboflow~\cite{roboflow}, plus
  auto-labeled frames from an earlier checkpoint; 70/20/10
  train/val/test split. Fine-tuned from COCO, 100 epochs,
  640$\times$640, augmentation: HSV jitter, 50\% horizontal flip,
  $\pm$0.3 scale, $\pm$0.1 translation, 50\% mosaic, $\pm$5$^\circ$
  rotation. Runs at 320$\times$320 on the left camera at inference;
  backbone for 4.6--4.7. [TODO: accuracy (mAP)]
  \item \textbf{4.6 Angular Displacement}: HSV segmentation on the
  YOLOv8-detected cylinder crop: blue/white masks, area-weighted
  centroids (all connected regions above a pixel threshold, not just
  the largest, so a split blob doesn't skew it), angle $=
  \operatorname{atan2}$ between them ($0^\circ$ = upright). Feeds two
  independent checks at placement: upright ($|\theta| \le 20^\circ$)
  and correct color down (vs.\ the cued target); a cylinder can pass
  one and fail the other. Fires 0.5~s after target hit, up to 3
  retries.
  \item \textbf{4.7 Object State Detection}: Upright/color-down
  scoring from 4.6 against the lit target peg. Separately, a 4-class
  state classifier (stationary/held/dropped/lost, out of the ECM's
  view; single cylinder only). First a rule-based version
  (pixel displacement/velocity, spike/settle counters); not accurate
  enough (held vs.\ dropped), so its own output bootstrapped training
  data (state, pixel velocity, bbox position, PSM1 pose+quaternion),
  hand-corrected against a video recording as ground truth. XGBoost
  classifier~\cite{xgboost} trained on that (300 trees, depth 6, LR
  0.05; features include a 4-frame lag window + deltas); validated
  with grouped stratified $k$-fold CV, scored by macro-F1 (weights
  rare states like dropped/lost equally with common ones). [TODO:
  F1-macro mean $\pm$ std]
  \item \textbf{4.8 3D Localization (Post-Hoc)}: One cylinder at a
  time. Early attempt reused the YOLOv8 approach (4.5) on the depth
  camera's own view (record, label on Roboflow~\cite{roboflow}, train,
  run live); not good enough (lag, low accuracy, poorly localized
  centroid). Rebuilt on SAM2~\cite{sam2}, post-hoc, five stages: (1)
  exposure tuning, hand-tuned exposure/gain/laser/white balance
  against a live view until all pegs are clearly visible; (2) peg
  registration, up axis from 3 equal-height clicks, then all 12 peg
  tops clicked and averaged over 30 frames each, reviewed in 3D for
  outliers; (3) record \& extract, color+depth recorded to one file
  per session (PSM kinematics recorded separately), then frames +
  aligned depth extracted; (4) segmentation, SAM2 segments the
  cylinder from a few clicked points, propagated through the clip;
  (5) visualization, each frame's masked pixels deprojected to 3D,
  median = cylinder position, compared to the calibrated peg
  positions for distance/height, giving a 3D trajectory over the
  board. SAM2 propagation runs at $\sim$5~fps, well under real-time,
  so this stays post-hoc; serves as an offline validation reference
  for path length (4.3),
  not a delivered real-time measure.
\end{itemize}


\begin{table}[h]
\centering
\resizebox{\textwidth}{!}{%
\begin{tabular}{|p{5.2cm}|p{3.2cm}|p{5.6cm}|p{5.6cm}|p{2.4cm}|}
\hline
\textbf{Data / Measure} & \textbf{Timing} & \textbf{Trigger} &
\textbf{Feedback} & \textbf{Cylinders} \\
\hline
Trial time &
Real-time &
Continuous, lift-to-placement &
Progress bar (GUI) &
Multiple \\[4pt]
\hline
Path length &
Real-time &
Continuous while clock runs &
Numeric running total &
Multiple \\[4pt]
\hline
Rate of orientation change &
Real-time &
Continuous while clock runs &
Numeric running value &
Multiple \\[4pt]
\hline
Object detection (YOLOv8~\cite{yolov8}) &
Real-time, delayed &
Continuous, event-triggered at placement &
Backbone for orientation \& state below &
Multiple \\[4pt]
\hline
Cylinder orientation (HSV) &
Real-time, delayed &
Continuous, event-triggered at placement &
Two binary indicators (upright, correct color down) &
Multiple \\[4pt]
\hline
Object state (4 classes) &
Real-time &
Continuous &
Status indicator &
Single \\[4pt]
\hline
Contact force &
Real-time, delayed &
Continuous, read live; shown after trial ends &
Popup: peak force, threshold crossings, force/time graph &
Multiple \\[4pt]
\hline
Additional camera &
Real-time &
Continuous, pedal-triggered &
Additional camera view &
Multiple \\[4pt]
\hline
3D localization (SAM2~\cite{sam2}) &
Post-hoc &
Offline processing &
Not shown to trainee; offline plot &
Single \\
\hline
\end{tabular}%
}
\caption{Data collected by the pipeline, its timing class, trigger, the
feedback form it reaches the trainee in, and whether the technique
supports single or multiple cylinders on the task pad at once, matching
\texttt{tab:data\_timing} in \texttt{sections/methods.tex}.}
\label{tab:outline_data_timing}
\end{table}

\newpage

\subsection*{5. Results \& Evaluation: PLACEHOLDER}
% Drafted in the thesis (results.tex); numbers still pending.
Bench trials run by the author, not a trainee study (that is future
work, Section 7). Organized around the merged pipeline rather than
one subsection per measure, since the point is what one trial run
through the whole system produces together, with timing broken out
for a direct, quantitative check.
\begin{itemize}
  \item \textbf{5.1 An Integrated Trial}: 1--3 representative trials
  (clean success, a drop, a wrong-color placement) walked through end
  to end, in pipeline order: trial time/path length (GUI), orientation
  rate (terminal), upright/color-down + object-state at placement,
  force popup after. \emph{[Figure/screenshot: GUI + force popup per
  trial]} [TODO: run and log the trials]
  \item \textbf{5.2 Timing}: Measured, not assumed, latency per
  module: GUI/kinematics update rate vs.\ the 200~Hz assumption;
  orientation-check delay (0.5~s settle + processing + retries) vs.\
  measured; object-state classifier per-frame inference time; force
  popup delay after trial end. [TODO: all numbers; may be computable
  from existing bench-trial logs if a completion timestamp is added]
\end{itemize}


\subsection*{6. Discussion: PLACEHOLDER}
% Drafted in the thesis (discussion.tex); two pieces pending Results.
\begin{itemize}
  \item \textbf{6.1 Does This Close the Gap?}: Revisits Section 2's
  three-way gap (automated / real-time / physical hardware); argues
  this system closes it, with 3D localization's post-hoc timing as
  the one exception, not itself a gap since it was never delivered as
  a real-time measure.
  \item \textbf{6.2 What the Timing Means}: Interprets the 5.2 latency
  numbers, e.g.\ against human reaction time, and against Laca et
  al.'s real-time-loop/non-real-time-measurement split (Section 2).
  [TODO: pending 5.2 numbers]
  \item \textbf{6.3 Limitations}: Bench trials only, no trainee study
  yet; object-state classifier reliability not yet established
  (rule-based labeler struggled at held/dropped, XGBoost accuracy
  pending); 3D localization post-hoc only, real-time attempt
  abandoned; object-state \& 3D localization both single-cylinder
  only, despite the task allowing multiple; HSV thresholds tuned for
  one camera/lighting setup, untested elsewhere.
  \item \textbf{6.4 What This Sets Up}: Bridges to future work
  (Section 7): the redesigned task board and closed-loop adaptive
  task from the formative study (3.2), plus a trainee study of the
  system itself.
\end{itemize}


\subsection*{7. Conclusion \& Future Work}
\begin{itemize}
  \item Summary of contribution (real-time, automated, on-robot,
  no human rater, one task).
  \item Outlook: closed-loop, difficulty-adaptive task (farther peg,
  tighter orientation tolerance, difficulty set by prior trial score).
\end{itemize}

\bibliographystyle{plain}
\bibliography{references}

\end{document}
